% GENERATED FILE. Edit canonical lessons, then run npm run generate:books.
% canonical-source-hash: fnv1a64:8e77d877fe6d115a
% canonical-lessons: KA-C529-vayaska, KA-C529-hadihareya, KA-C529-ekanta, KA-C529-moha, KA-C529-baduku

\chapter{Stages of Life}
\label{ch:ka-stages-of-life}
\hlchaptermodality{529}

\begin{chapteropening}
\textbf{By the end of this chapter:} \emph{I can say an adult, the teenage years, solitude, infatuation and to live.}

Complete the last lesson of chapter 529: I can say an adult, the teenage years, solitude, infatuation and to live.
\end{chapteropening}

\section[\texorpdfstring{\kn{ವಯಸ್ಕ} (vayaska)}{ವಯಸ್ಕ (vayaska)}]{\kn{ವಯಸ್ಕ} (vayaska) — an adult}

\label{lesson:KA-C529-vayaska}



\begin{quote}
Before the new one: say the Kannada for dear, then the Kannada for to make up.
\end{quote}

\subsection*{You'll want to know: \kn{ವಯಸ್ಕ}}
\textbf{\kn{ವಯಸ್ಕ}} — \emph{vayaska} — \textquotedblleft{}an adult\textquotedblright{}.

\begin{quote}
\kn{ಹದಿನೆಂಟು} \kn{ವರ್ಷದ} \kn{ನಂತರ} \kn{ಮಕ್ಕಳು} \kn{ವಯಸ್ಕರು} \kn{ಆಗುತ್ತಾರೆ}.
\end{quote}

\begin{quote}
\emph{hadineṇṭu varṣada nantara makkaḷu vayaskaru āguttāre.}
\end{quote}

\begin{quote}
After eighteen, children become adults.
\end{quote}

Plural: \kn{ವಯಸ್ಕರು}, \emph{vayaskaru}.

Built on \kn{ವಯಸ್ಸು}, age.

\begin{cousinweb}
\textbf{\kn{ವಯಸ್ಕ}} (\emph{vayaska}) is built on Sanskrit \emph{vayas}, \textquotedblleft{}age, vigour of life\textquotedblright{}, which Monier-Williams also lists in compounds with -\emph{vayaska} (Monier-Williams).
\end{cousinweb}

\begin{grammarlens}[title={What you've built}]
One more word for describing what you lived and felt.
\end{grammarlens}

\subsection*{Your turn}
\begin{itemize}
\raggedright
  \item \emph{Say it:} \emph{vayaska}
  \item \emph{Say it:} \emph{vayaska}, once more
  \item \emph{Say it:} \emph{hadineṇṭu varṣada nantara makkaḷu vayaskaru āguttāre.} — After eighteen, children become adults.
  \item \emph{Recall:} say the Kannada for dear, then the Kannada for to make up, then say \emph{vayaska} again
\end{itemize}

\begin{tcolorbox}[breakable,colback=teal!4,colframe=teal!35!black,title={Before you move on}]
What does \kn{ವಯಸ್ಕ} mean? (\textquotedblleft{}An adult\textquotedblright{}.) Say it once more.
\end{tcolorbox}
\section[\texorpdfstring{\kn{ಹದಿಹರೆಯ} (hadihareya)}{ಹದಿಹರೆಯ (hadihareya)}]{\kn{ಹದಿಹರೆಯ} (hadihareya) — the teenage years}

\label{lesson:KA-C529-hadihareya}



\begin{quote}
Before the new one: say the Kannada for to make up, then the Kannada for an adult.
\end{quote}

\subsection*{You'll want to know: \kn{ಹದಿಹರೆಯ}}
\textbf{\kn{ಹದಿಹರೆಯ}} — \emph{hadihareya} — \textquotedblleft{}the teenage years\textquotedblright{}.

\begin{quote}
\kn{ಹದಿಹರೆಯದಲ್ಲಿ} \kn{ಮಕ್ಕಳ} \kn{ಸ್ವಭಾವ} \kn{ಬದಲಾಗುತ್ತದೆ}.
\end{quote}

\begin{quote}
\emph{hadihareyadalli makkaḷa svabhāva badalāguttade.}
\end{quote}

\begin{quote}
Children's nature changes in the teenage years.
\end{quote}

\kn{ಹದಿ}- is the \textquotedblleft{}-teen\textquotedblright{} of \kn{ಹದಿನಾರು}, sixteen; \kn{ಹರೆಯ} is youth.

\begin{cousinweb}
\textbf{\kn{ಹದಿಹರೆಯ}} (\emph{hadihareya}) is a Kannada compound of \emph{hadi-}, the \textquotedblleft{}ten\textquotedblright{} in \kn{ಹದಿನಾರು} (DEDR 3918 lists Kannada \emph{padi-}, ten, in compounds), and \emph{hareya}, \textquotedblleft{}youth\textquotedblright{}; the origin of \emph{hareya} is not settled.
\end{cousinweb}

\begin{grammarlens}[title={What you've built}]
One more word for describing what you lived and felt.
\end{grammarlens}

\subsection*{Your turn}
\begin{itemize}
\raggedright
  \item \emph{Say it:} \emph{hadihareya}
  \item \emph{Say it:} \emph{hadihareya}, once more
  \item \emph{Say it:} \emph{hadihareyadalli makkaḷa svabhāva badalāguttade.} — Children's nature changes in the teenage years.
  \item \emph{Recall:} say the Kannada for to make up, then the Kannada for an adult, then say \emph{hadihareya} again
\end{itemize}

\begin{tcolorbox}[breakable,colback=teal!4,colframe=teal!35!black,title={Before you move on}]
What does \kn{ಹದಿಹರೆಯ} mean? (\textquotedblleft{}The teenage years\textquotedblright{}.) Say it once more.
\end{tcolorbox}
\section[\texorpdfstring{\kn{ಏಕಾಂತ} (ēkānta)}{ಏಕಾಂತ (ēkānta)}]{\kn{ಏಕಾಂತ} (ēkānta) — solitude}

\label{lesson:KA-C529-ekanta}



\begin{quote}
Before the new one: say the Kannada for an adult, then the Kannada for the teenage years.
\end{quote}

\subsection*{You'll want to know: \kn{ಏಕಾಂತ}}
\textbf{\kn{ಏಕಾಂತ}} — \emph{ēkānta} — \textquotedblleft{}solitude\textquotedblright{}.

\begin{quote}
\kn{ವಯಸ್ಸಾದ} \kn{ಮೇಲೆ} \kn{ಅಜ್ಜನಿಗೆ} \kn{ಏಕಾಂತ} \kn{ಇಷ್ಟ}.
\end{quote}

\begin{quote}
\emph{vayassāda mēle ajjanige ēkānta iṣṭa.}
\end{quote}

\begin{quote}
In old age grandfather likes solitude.
\end{quote}

Chosen quiet; \kn{ಒಂಟಿತನ} is loneliness you did not choose.

\begin{cousinweb}
\textbf{\kn{ಏಕಾಂತ}} (\emph{ēkānta}) is borrowed from Sanskrit \emph{ekānta}, \textquotedblleft{}a lonely or retired place\textquotedblright{} (Monier-Williams).
\end{cousinweb}

\begin{grammarlens}[title={What you've built}]
One more word for describing what you lived and felt.
\end{grammarlens}

\subsection*{Your turn}
\begin{itemize}
\raggedright
  \item \emph{Say it:} \emph{ēkānta}
  \item \emph{Say it:} \emph{ēkānta}, once more
  \item \emph{Say it:} \emph{vayassāda mēle ajjanige ēkānta iṣṭa.} — In old age grandfather likes solitude.
  \item \emph{Recall:} say the Kannada for an adult, then the Kannada for the teenage years, then say \emph{ēkānta} again
\end{itemize}

\begin{tcolorbox}[breakable,colback=teal!4,colframe=teal!35!black,title={Before you move on}]
What does \kn{ಏಕಾಂತ} mean? (\textquotedblleft{}Solitude\textquotedblright{}.) Say it once more.
\end{tcolorbox}
\section[\texorpdfstring{\kn{ಮೋಹ} (mōha)}{ಮೋಹ (mōha)}]{\kn{ಮೋಹ} (mōha) — infatuation; enchantment}

\label{lesson:KA-C529-moha}



\begin{quote}
Before the new one: say the Kannada for the teenage years, then the Kannada for solitude.
\end{quote}

\subsection*{You'll want to know: \kn{ಮೋಹ}}
\textbf{\kn{ಮೋಹ}} — \emph{mōha} — \textquotedblleft{}infatuation; enchantment\textquotedblright{}.

\begin{quote}
\kn{ಹದಿಹರೆಯದಲ್ಲಿ} \kn{ಅವನಿಗೆ} \kn{ಸಿನಿಮಾದ} \kn{ಮೇಲೆ} \kn{ಮೋಹ} \kn{ಇತ್ತು}.
\end{quote}

\begin{quote}
\emph{hadihareyadalli avanige sinimāda mēle mōha ittu.}
\end{quote}

\begin{quote}
In his teens he was infatuated with the cinema.
\end{quote}

For what: … \kn{ಮೇಲೆ} \kn{ಮೋಹ}.

\begin{cousinweb}
\textbf{\kn{ಮೋಹ}} (\emph{mōha}) is borrowed from Sanskrit \emph{moha}, from \emph{muh}, \textquotedblleft{}to be bewildered\textquotedblright{}: \textquotedblleft{}loss of consciousness, bewilderment, infatuation\textquotedblright{} (Monier-Williams).
\end{cousinweb}

\begin{grammarlens}[title={What you've built}]
One more word for describing what you lived and felt.
\end{grammarlens}

\subsection*{Your turn}
\begin{itemize}
\raggedright
  \item \emph{Say it:} \emph{mōha}
  \item \emph{Say it:} \emph{mōha}, once more
  \item \emph{Say it:} \emph{hadihareyadalli avanige sinimāda mēle mōha ittu.} — In his teens he was infatuated with the cinema.
  \item \emph{Recall:} say the Kannada for the teenage years, then the Kannada for solitude, then say \emph{mōha} again
\end{itemize}

\begin{tcolorbox}[breakable,colback=teal!4,colframe=teal!35!black,title={Before you move on}]
What does \kn{ಮೋಹ} mean? (\textquotedblleft{}Infatuation; enchantment\textquotedblright{}.) Say it once more.
\end{tcolorbox}
\section[\texorpdfstring{\kn{ಬದುಕು} (baduku)}{ಬದುಕು (baduku)}]{\kn{ಬದುಕು} (baduku) — to live; to survive}

\label{lesson:KA-C529-baduku}



\begin{quote}
Before the new one: say the Kannada for solitude, then the Kannada for infatuation.
\end{quote}

\subsection*{You'll want to know: \kn{ಬದುಕು}}
\textbf{\kn{ಬದುಕು}} — \emph{baduku} — \textquotedblleft{}to live; to survive\textquotedblright{}.

\begin{quote}
\kn{ಅಜ್ಜ} \kn{ತೊಂಬತ್ತು} \kn{ವರ್ಷ} \kn{ಬದುಕಿದರು}.
\end{quote}

\begin{quote}
\emph{ajja tombattu varṣa badukidaru.}
\end{quote}

\begin{quote}
Grandfather lived to ninety.
\end{quote}

Past (I …): \kn{ಬದುಕಿದೆ}, \emph{badukide}.

Also a noun, life: \kn{ಬದುಕು} \kn{ಕಷ್ಟ}, life is hard.

\begin{cousinweb}
\textbf{\kn{ಬದುಕು}} (\emph{baduku}) is native Dravidian: DEDR 5372 lists Kannada \emph{barduku} and \emph{baduku}, \textquotedblleft{}to live, be alive, make a livelihood\textquotedblright{}.
\end{cousinweb}

\begin{grammarlens}[title={What you've built}]
One more word for describing what you lived and felt.
\end{grammarlens}

\subsection*{Your turn}
\begin{itemize}
\raggedright
  \item \emph{Say it:} \emph{baduku}
  \item \emph{Say it:} \emph{baduku}, once more
  \item \emph{Say it:} \emph{ajja tombattu varṣa badukidaru.} — Grandfather lived to ninety.
  \item \emph{Recall:} say the Kannada for solitude, then the Kannada for infatuation, then say \emph{baduku} again
\end{itemize}

\begin{tcolorbox}[breakable,colback=teal!4,colframe=teal!35!black,title={Before you move on}]
What does \kn{ಬದುಕು} mean? (\textquotedblleft{}To live; to survive\textquotedblright{}.) Say it once more.
\end{tcolorbox}
